\documentclass[aspectratio=169,11pt]{beamer}
\usepackage{paperdeck}
\title[IL-TUR]{IL-TUR: Benchmark for Indian Legal Text Understanding and Reasoning}
\subtitle{Joshi, Paul, Sharma, Goyal, Ghosh and Modi (ACL 2024), arXiv:2407.05399}
\author[Paper summary]{Paper summary for the CS-781 capstone}
\date{}
\begin{document}

\begin{frame}\titlepage\end{frame}

\paperglance
{A.\ Joshi, S.\ Paul, A.\ Sharma, P.\ Goyal, S.\ Ghosh, A.\ Modi. \emph{IL-TUR: Benchmark for Indian Legal Text Understanding and Reasoning}. ACL 2024 (arXiv v2, November 2024). IIT Kanpur and IIT Kharagpur.}
{A single benchmark of eight Indian legal NLP tasks, with unified data formats, baselines (including LLMs) and a public leaderboard. Current models are far from solving it.}
{This is the capstone's \textbf{base paper}. Phases 2--7 are IL-TUR tasks: L-NER, RR, CJPE, LSI, SUMM and PCR.}
{Main deck: ``Capstone at a glance'' and every Phase 2--7 slide.}

\begin{frame}{Roadmap}
  \begin{multicols}{2}\small\tableofcontents\end{multicols}
\end{frame}

% ------------------------------------------------------------------
\section{The problem}

\begin{frame}{Why legal NLP for India?}
  \begin{columns}[T]
    \begin{column}{0.5\textwidth}
      \small
      \begin{itemize}
        \item India has about \textbf{50 million pending cases} across district, state and federal courts (National Judicial Data Grid, 2023).
        \item ``Justice delayed is justice denied.''
        \item Legal work runs on documents, so NLP could help lawyers and judges read, search and summarise faster.
        \item But legal NLP tasks existed in isolation, with different formats and metrics, so progress was hard to track.
      \end{itemize}
    \end{column}
    \begin{column}{0.46\textwidth}
      \begin{exampleblock}{The idea in one line}
        \small
        Do for Indian legal NLP what GLUE did for general NLP: one benchmark, one format, one leaderboard.
      \end{exampleblock}
      \begin{alertblock}{Not meant to replace judges}
        \small
        The goal is to \emph{augment} legal practitioners, e.g.\ by suggesting statutes, surfacing precedents or summarising, not to automate decisions.
      \end{alertblock}
    \end{column}
  \end{columns}
\end{frame}

\begin{frame}{Why a separate benchmark for legal text?}
  \small
  \begin{tabularx}{\textwidth}{@{}lX@{}}
    \toprule
    Challenge & What it means, with an example\\
    \midrule
    1.\ Domain vocabulary & Everyday words take special meanings: ``consideration'', ``relief'', ``ratio'', ``stay''\\
    2.\ Very long documents & An Indian Supreme Court judgment averages about \textbf{4,000 words}; BERT reads only 512 tokens\\
    3.\ Unstructured and noisy & Many documents are typed manually in courts, with typos and no consistent sections\\
    4.\ Sub-domains differ & A model trained on civil law may fail on banking or insurance law: domain adaptation is hard\\
    5.\ Explainability required & A lawyer or judge must see \emph{why} a model predicts an outcome\\
    \bottomrule
  \end{tabularx}
  \vspace{2mm}
  \begin{exampleblock}{Mixed legal system}
    \small
    India is mainly \textbf{common law}: higher courts can overrule precedents, so outcomes need reasoning beyond the statute book. Some areas (e.g.\ banking, insurance) follow \textbf{civil law}. IL-TUR covers both.
  \end{exampleblock}
\end{frame}

\begin{frame}{How IL-TUR differs from other legal benchmarks}
  \centering
  \footnotesize
  \begin{tabular}{@{}llll@{}}
    \toprule
    Benchmark & Jurisdiction & Task types & Languages\\
    \midrule
    LexGLUE & US, EU & classification & English\\
    LEXTREME & EU, Brazil & classification & 24 EU\\
    FAIRLEX & EU, US, China, CH & fairness & several\\
    LBOX & Korea & class., generation & Korean\\
    LegalBench & mostly US & class., generation & English\\
    LawBench & China & class., gen., extraction & Chinese\\
    \textbf{IL-TUR} & \textbf{India} & \textbf{class., retrieval, gen., extraction} & \textbf{English + 9 Indian}\\
    \bottomrule
  \end{tabular}
  \vspace{2mm}

  \raggedright
  \begin{exampleblock}{Key differences}
    \small
    Tasks are grounded in the \textbf{real legal workflow}, not exam-style questions. Documents are \textbf{long}: about 4,000 words, against about 200 in LegalBench. Tasks include retrieval, generation and explanation, and both common and civil law are covered.
  \end{exampleblock}
\end{frame}

% ------------------------------------------------------------------
\section{The benchmark}

\begin{frame}{Design principles}
  \begin{columns}[T]
    \begin{column}{0.5\textwidth}
      \small
      \begin{enumerate}
        \item \textbf{Legal understanding and world knowledge}: tasks need real legal knowledge plus common sense.
        \item \textbf{Difficulty}: not solvable by a layperson; hard for current models.
        \item \textbf{Language}: English and the most frequent Indian court languages.
        \item \textbf{Automatic evaluation} with metrics that align with human judgement.
        \item \textbf{Public availability}: CC BY-NC-SA data and a public leaderboard.
      \end{enumerate}
    \end{column}
    \begin{column}{0.46\textwidth}
      \begin{block}{Harmonisation}
        \small
        Several tasks existed before, each with its own format and metric. IL-TUR converts all of them to a \textbf{uniform JSON format}, collects training scripts, and defines a standard evaluation.
      \end{block}
      \begin{exampleblock}{Anonymisation}
        \small
        Named entities are anonymised in RR, CJPE, BAIL, LSI and PCR data, to address ethical concerns and to stop models from learning biases tied to names.
      \end{exampleblock}
    \end{column}
  \end{columns}
\end{frame}

\begin{frame}{The eight tasks}
  \footnotesize
  \begin{tabularx}{\textwidth}{@{}llrlX@{}}
    \toprule
    Task & Dataset & Avg words & Type & Key skill\\
    \midrule
    L-NER & 105 docs, 650k words & 6,180 & sequence labelling & legal roles of entities (new task)\\
    RR & 21,184 sentences & 25,796 & multi-class & legal semantics per sentence\\
    CJPE & ILDC, 34k docs & 3,336 & class.\ + extraction & understanding and reasoning, with explanation\\
    BAIL & HLDC, 176k docs (Hindi) & 86 & classification & reasoning in Hindi\\
    LSI & ILSI, 65k samples & 2,406 & multi-label & applying statutes to facts\\
    PCR & IL-PCR, 7,070 docs & 8,096 & retrieval & legal relevance between cases\\
    SUMM & In-Abs, 7,130 docs & 4,376 & generation & legal understanding and writing\\
    L-MT & MILPaC, 17,853 pairs & 49 & generation & English $\leftrightarrow$ 9 Indian languages\\
    \bottomrule
  \end{tabularx}
  \vspace{2mm}

  \small
  Note the lengths: prior case retrieval documents average about 8,100 words, far beyond a 1,024-token window.
\end{frame}

\begin{frame}{Tasks in the litigation process, and in our capstone}
  \begin{columns}[T]
    \begin{column}{0.5\textwidth}
      \centering
      \begin{tikzpicture}[node distance=2.5mm]
        \node[box,minimum width=46mm] (f) {read facts: \textbf{L-NER}, \textbf{RR}, \textbf{SUMM}};
        \node[hot,below=of f,minimum width=46mm] (s) {find applicable law: \textbf{LSI}};
        \node[hot,below=of s,minimum width=46mm] (p) {find precedents: \textbf{PCR}};
        \node[hot,below=of p,minimum width=46mm] (d) {decide: \textbf{CJPE}, \textbf{BAIL}};
        \node[box,below=of d,minimum width=46mm] (t) {re-appeal: translate, \textbf{L-MT}};
        \draw[arr] (f)--(s); \draw[arr] (s)--(p); \draw[arr] (p)--(d); \draw[arr] (d)--(t);
      \end{tikzpicture}
    \end{column}
    \begin{column}{0.46\textwidth}
      \footnotesize
      \begin{tabular}{@{}ll@{}}
        \toprule
        Capstone phase & IL-TUR task\\
        \midrule
        Phase 2 & L-NER\\
        Phase 3 & RR (CL and IT domains)\\
        Phase 4 & CJPE (prediction part)\\
        Phase 5 & LSI\\
        Phase 6 & SUMM\\
        Phase 7 & PCR\\
        \bottomrule
      \end{tabular}
      \vspace{1mm}

      \small
      BAIL (Hindi) and L-MT (translation) are not capstone phases. Phase 1, the foundation model, is the capstone's own addition.
    \end{column}
  \end{columns}
\end{frame}

\begin{frame}{Easy examples of each capstone task}
  \scriptsize
  \begin{tabularx}{\textwidth}{@{}lXl@{}}
    \toprule
    Task & Input (simplified) & Output\\
    \midrule
    L-NER & ``Appeal by \emph{Ram Kumar} against the order of the \emph{High Court of Delhi} dated \emph{3 May 2019}\ldots'' & APP, COURT, DATE spans\\
    RR & ``The accused was arrested on 3 May.'' / ``Whether Section 302 applies.'' & Fact / Issue\\
    CJPE & full judgment with the final order removed & 1 (appeal accepted)\\
    LSI & facts: ``the accused, together with two others, planned and caused the death of\ldots'' & IPC 302, IPC 120B\\
    PCR & query judgment with citations removed, and a pool of candidate judgments & top-7 relevant cases\\
    SUMM & 4,000-word judgment & a few-paragraph summary\\
    \bottomrule
  \end{tabularx}
  \vspace{2mm}
  \begin{exampleblock}{Why L-NER is not ordinary NER}
    \small
    A standard NER model tags the judge, petitioner and respondent all as PERSON. L-NER must know their \emph{legal roles}, which needs legal understanding. Its gold data was annotated by law students, with entity definitions from legal academics.
  \end{exampleblock}
\end{frame}

% ------------------------------------------------------------------
\section{Results}

\begin{frame}{Best baselines per task}
  \centering
  \footnotesize
  \begin{tabular}{@{}llll@{}}
    \toprule
    Task & Best score & Metric & Model\\
    \midrule
    L-NER & 48.58\% & strict macro-F1 & InLegalBERT + CRF\\
    RR & 69.01\% & macro-F1 & MTL-BERT\\
    CJPE & 81.31\% & macro-F1 & InLegalBERT + BiLSTM\\
    BAIL & 81\% & macro-F1 & TF-IDF + IndicBERT\\
    LSI & 28.08\% & macro-F1 & LeSICiN (graph-based)\\
    PCR & 39.15\% & $\mu$F1@K & event-based model\\
    SUMM & 0.33 / 0.86 & ROUGE-L / BERTScore & Legal-LED\\
    L-MT & 0.28 / 0.32 / 0.57 & BLEU / GLEU / chrF++ & Microsoft Translator\\
    \bottomrule
  \end{tabular}
  \vspace{2mm}

  \raggedright
  \begin{exampleblock}{Reading it}
    \small
    Strict L-NER is below 50\% (exact boundaries are hard), LSI below 30\% macro-F1 (rare statutes), and PCR below 40\% (legal relevance is subtle).
    These are the scores our Phase 2--7 systems will be compared against.
  \end{exampleblock}
\end{frame}

\begin{frame}{Does in-domain pre-training help?}
  \centering
  \footnotesize
  \begin{tabular}{@{}llccc@{}}
    \toprule
    Task & Setup & BERT & LegalBERT & InLegalBERT\\
    \midrule
    L-NER & flat, chunks & 39.59 & 45.58 & \textbf{48.58}\\
    RR & flat, sentences & 58 & 54 & 58\\
    CJPE & hierarchical & 71.14 & 78.21 & \textbf{81.31}\\
    LSI & hierarchical & 18.44 & 21.74 & \textbf{26.23}\\
    PCR & flat, chunks & \textbf{9.24} & 8.67 & 7.57\\
    \bottomrule
  \end{tabular}
  \vspace{1mm}

  {\scriptsize Macro-F1 (strict for L-NER). LegalBERT: pre-trained on EU/US legal text. InLegalBERT: on Indian legal text.}
  \vspace{2mm}

  \raggedright
  \begin{exampleblock}{Lesson for our Phase 1}
    \small
    Except for PCR, performance rises from general BERT, to foreign legal BERT, to \textbf{Indian} legal BERT: it tracks how in-domain the pre-training was.
    This supports our plan to pre-train from scratch on Indian legal text. For CJPE, \emph{hierarchical} models (encoder + BiLSTM over chunks) worked best, because judgments exceed 512 tokens.
  \end{exampleblock}
\end{frame}

\begin{frame}{Large language models are far behind}
  \begin{columns}[T]
    \begin{column}{0.56\textwidth}
      \footnotesize
      Best GPT-3.5 / GPT-4 score over 0-, 1- and 2-shot prompts, vs the task-specific best:
      \vspace{1mm}

      \begin{tabular}{@{}lccc@{}}
        \toprule
        Task (metric) & GPT-3.5 & GPT-4 & Best baseline\\
        \midrule
        L-NER (strict mF1) & 32.84 & 24.03 & \textbf{48.58}\\
        RR (mF1) & 30.95 & 38.18 & \textbf{69.01}\\
        CJPE (mF1) & 56.74 & 68.29 & \textbf{81.31}\\
        BAIL (mF1) & 61.0 & 66.67 & \textbf{81}\\
        LSI (mF1) & 22.61 & 23.99 & \textbf{28.08}\\
        SUMM (ROUGE-L) & 0.22 & 0.23 & \textbf{0.33}\\
        L-MT (BLEU) & 0.26 & \textbf{0.36} & 0.28\\
        \bottomrule
      \end{tabular}
    \end{column}
    \begin{column}{0.4\textwidth}
      \begin{exampleblock}{Reading it}
        \small
        \begin{itemize}
          \item Task-specific models beat GPT-4 on every task except translation.
          \item Long contexts and structured outputs (spans, one label per sentence) are hard to get from prompting.
          \item Extra examples did not reliably help; one-shot was sometimes worse than zero-shot.
          \item PCR could not be tried: the candidate pool exceeds even GPT-4's context.
        \end{itemize}
      \end{exampleblock}
    \end{column}
  \end{columns}
\end{frame}

\begin{frame}{Generalising across jurisdictions}
  \small
  Most tasks are jurisdiction-specific (statutes, bail practice), but summarisation can be tested across countries: train on Indian \textbf{In-Abs}, then test zero-shot on UK Supreme Court summaries (\textbf{UK-Abs}).
  \vspace{2mm}

  \centering
  \footnotesize
  \begin{tabular}{@{}llccc@{}}
    \toprule
    Model & Trained on & ROUGE-1 & ROUGE-L & BERTScore\\
    \midrule
    SummaRuNNer & In-Abs & \textbf{0.604} & 0.225 & 0.828\\
    & UK-Abs & 0.493 & \textbf{0.274} & \textbf{0.849}\\
    Legal-LED & In-Abs & \textbf{0.557} & 0.242 & 0.844\\
    & UK-Abs & 0.482 & \textbf{0.264} & \textbf{0.851}\\
    \bottomrule
  \end{tabular}
  \vspace{2mm}

  \raggedright
  \begin{exampleblock}{Reading it}
    \small
    Models trained only on Indian data summarise UK judgments decently: higher ROUGE-1, and comparable ROUGE-L and BERTScore. Cross-jurisdiction adaptation is left as future work.
  \end{exampleblock}
\end{frame}

% ------------------------------------------------------------------
\section{Assessment and outlook}

\begin{frame}{Strengths and limitations}
  \begin{columns}[T]
    \begin{column}{0.48\textwidth}
      \begin{exampleblock}{Strengths}
        \small
        \begin{itemize}
          \item First unified benchmark for Indian legal NLP.
          \item Realistic, workflow-grounded tasks with long documents.
          \item Several new tasks and datasets (e.g.\ L-NER).
          \item Uniform formats, public data and leaderboard.
          \item Strong baselines, including LLM comparisons.
          \item Anonymisation to reduce bias.
        \end{itemize}
      \end{exampleblock}
    \end{column}
    \begin{column}{0.48\textwidth}
      \begin{alertblock}{Limitations (mostly stated by the authors)}
        \small
        \begin{itemize}
          \item Not every legal sub-domain is covered in each task.
          \item Multilingual only for L-MT; BAIL only in Hindi.
          \item Some datasets were built with ML tools (e.g.\ fact extraction for LSI), so labels may be imperfect.
          \item Explainability is addressed mainly in CJPE.
          \item LLMs were evaluated on subsets, because of cost.
          \item Only eight tasks so far.
        \end{itemize}
      \end{alertblock}
    \end{column}
  \end{columns}
\end{frame}

\begin{frame}{Future research directions}
  \begin{columns}[T]
    \begin{column}{0.5\textwidth}
      \small
      \textbf{Stated by the authors}
      \begin{itemize}
        \item Add tasks, e.g.\ \textbf{legal coreference resolution} for information extraction and knowledge graphs.
        \item A \textbf{single model for all tasks}, e.g.\ mT5 or multi-task adapters.
        \item More Indian languages and lower-court data.
        \item Refined datasets built with legal professionals.
        \item Cross-jurisdiction domain adaptation.
        \item Two-stage retrieval so LLMs can attempt PCR.
      \end{itemize}
    \end{column}
    \begin{column}{0.46\textwidth}
      \begin{block}{Directions our capstone explores}
        \small
        \begin{itemize}
          \item A \textbf{self-hosted foundation model} trained only on Indian legal text and reused across six tasks, close to the authors' ``one model for all tasks'' goal.
          \item Parameter-efficient adaptation (LoRA, LORMA) and distillation for small, private deployments.
          \item Long-document strategies (chunking, hierarchies, ALiBi or RoPE) for CJPE, LSI and PCR.
        \end{itemize}
      \end{block}
    \end{column}
  \end{columns}
\end{frame}

\begin{frame}{Summary}
  \begin{itemize}
    \item India's court backlog motivates legal NLP; legal text is long, noisy, specialised and needs explanations.
    \item \textbf{IL-TUR}: eight tasks (L-NER, RR, CJPE, BAIL, LSI, PCR, SUMM, L-MT), English and 9 Indian languages, unified formats and a public leaderboard.
    \item Best baselines leave large gaps: strict L-NER 48.6\%, LSI 28.1\%, PCR 39.2\%.
    \item In-domain (Indian legal) pre-training helps on most tasks; hierarchical models help on long documents.
    \item GPT-3.5 and GPT-4 trail task-specific models on every task except translation.
    \item For us: Phases 2--7 \emph{are} IL-TUR tasks, and these baselines are the bar to beat.
  \end{itemize}
\end{frame}

\begin{frame}{Reference}
  \small
  Joshi, A., Paul, S., Sharma, A., Goyal, P., Ghosh, S., and Modi, A.\ (2024). IL-TUR: Benchmark for Indian Legal Text Understanding and Reasoning. In \emph{Proceedings of ACL 2024}. arXiv:2407.05399. Leaderboard: \url{https://exploration-lab.github.io/IL-TUR/}.
  \vspace{3mm}

  {\footnotesize Other work cited: Malik et al.\ (2021) ILDC; Malik et al.\ (2022) Semantic Segmentation of Legal Documents; Kapoor et al.\ (2022) HLDC; Paul et al.\ (2022) LeSICiN; Joshi et al.\ (2023) U-CREAT (event-based PCR); Paul et al.\ (2023) InLegalBERT; Chalkidis et al.\ (2022) LexGLUE; Guha et al.\ (2023) LegalBench.}
\end{frame}

\end{document}
